\documentclass{article}

\usepackage{amsmath,amssymb}
\usepackage{xurl}
\usepackage[hidelinks]{hyperref}
\setlength{\emergencystretch}{2em}

\input{command}

\title{Schmidt-number premise of Wolf's reduction criterion (Wolf eq.~(3.18))}
\author{}
\date{2026-06-24}

\begin{document}
\maketitle
\gapnote{scope-restriction}{resolved}

\section{Source statement}

For an integer $n\geq1$, define the reduction map
\begin{align}
    T_n(X)&=\tr(X)\Id-\frac{1}{n}X.
    \label{eq:wolf_reduction_criterion_schmidt_premise_reduction_map}
\end{align}
The reduction-criterion discussion in
\path{Notes/WolfNoteTexSource/ch03_positive_not_completely.tex}, lines
329--338, states that $T_n$ is $n$-positive and self-dual
\cite{Wolf2012Quantum}.

Let $\rho$ be a bipartite state, and let $\rho_1$ and $\rho_2$ be its reduced
density matrices on the first and second factors. Equation~(3.18) of
Ref.~\cite{Wolf2012Quantum} states that Schmidt number at most $n$ implies
\begin{align}
    n\rho_1\otimes\Id&\geq\rho,
    \label{eq:wolf_reduction_criterion_schmidt_premise_source_left}\\
    n\Id\otimes\rho_2&\geq\rho.
    \label{eq:wolf_reduction_criterion_schmidt_premise_source_right}
\end{align}
Writing $\operatorname{sn}(\rho)$ for the Schmidt number, the source assertion
is
\begin{align}
    \operatorname{sn}(\rho)\leq n
        &\Longrightarrow
        \left(n\rho_1\otimes\Id\geq\rho
            \ \text{and}\
            n\Id\otimes\rho_2\geq\rho\right).
    \label{eq:wolf_reduction_criterion_schmidt_premise_source_implication}
\end{align}

The source proof has two steps. A state of Schmidt number at most $n$ is a
convex combination of pure-state projectors whose vectors have Schmidt rank at
most $n$. The $n$-positivity of $T_n$ therefore gives
\begin{align}
    \operatorname{sn}(\rho)\leq n
        &\Longrightarrow
        (T_n\otimes\operatorname{id})(\rho)\geq0.
    \label{eq:wolf_reduction_criterion_schmidt_premise_schmidt_to_positivity}
\end{align}
The algebraic identity
\begin{align}
    (T_n\otimes\operatorname{id})(\rho)
        &=\Id\otimes\rho_2-\frac{1}{n}\rho
    \label{eq:wolf_reduction_criterion_schmidt_premise_tensor_identity}
\end{align}
then implies
\begin{align}
    (T_n\otimes\operatorname{id})(\rho)\geq0
        &\Longrightarrow
        n\Id\otimes\rho_2\geq\rho.
    \label{eq:wolf_reduction_criterion_schmidt_premise_operator_step}
\end{align}
Applying the same argument after swapping the tensor factors gives
(\ref{eq:wolf_reduction_criterion_schmidt_premise_source_left}).

\section{The initially formalized operator step}

The first formal theorems established only the implication in
(\ref{eq:wolf_reduction_criterion_schmidt_premise_operator_step}).\footnote{Formal
declarations:
    \leanid{Matrix.reductionCriterion_left} and
    \leanid{Matrix.reductionCriterion_right} in
    \doclink{QICLean/Channel/ReductionCriterion}.}
For a bipartite matrix $\rho$ and $n\geq1$, they prove
\begin{align}
    (T_n\otimes\operatorname{id})(\rho)\geq0
        &\Longrightarrow
        n\Id\otimes\rho_2\geq\rho,
    \label{eq:wolf_reduction_criterion_schmidt_premise_formal_left}\\
    (T_n\otimes\operatorname{id})(\rho^F)\geq0
        &\Longrightarrow
        n\rho_1\otimes\Id\geq\rho,
    \label{eq:wolf_reduction_criterion_schmidt_premise_formal_right}
\end{align}
where $\rho^F$ is obtained by swapping the tensor factors. The second
hypothesis is the swapped form of
$(\operatorname{id}\otimes T_n)(\rho)\geq0$. The identity
(\ref{eq:wolf_reduction_criterion_schmidt_premise_tensor_identity}) is also
formalized.\footnote{Formal declaration:
    \leanid{Matrix.tensorMapId_tEta_eq}.}

\section{Resolution of the Schmidt-number premise}

The formal predicate
\leanid{Matrix.HasSchmidtNumberLE} records that $\rho$ is a finite sum of
pure-state projectors whose vectors have Schmidt rank at most $n$.\footnote{The
declaration is in \doclink{QICLean/Channel/SchmidtNumber}.}
This is the convex decomposition used in Ref.~\cite{Wolf2012Quantum}, up to
normalization. At $n=1$, the predicate is equivalent to
separability.\footnote{Formal declaration:
    \leanid{Matrix.hasSchmidtNumberLE_one_iff_isSeparable}.}

For $1\leq n<D$, the pure-state theorem states that
\begin{align}
    \operatorname{SR}(\psi)\leq n
        &\Longrightarrow
        (T_n\otimes\operatorname{id})
        \!\left(\ketbra{\psi}{\psi}\right)\geq0.
    \label{eq:wolf_reduction_criterion_schmidt_premise_pure_state_step}
\end{align}
This is
\leanid{Matrix.tensorMapId_tEta_posSemidef_of_hasSchmidtRankLE}. The proof
writes $\psi$ using a maximally entangled vector and a square right-factor
matrix whose rank equals $\operatorname{SR}(\psi)$. It identifies the
ampliation with a right-factor compression of the Choi matrix and applies the
$n$-positivity of $T_n$.

Summing
(\ref{eq:wolf_reduction_criterion_schmidt_premise_pure_state_step}) over the
pure terms in a Schmidt-number decomposition proves
(\ref{eq:wolf_reduction_criterion_schmidt_premise_schmidt_to_positivity}).\footnote{Formal
declaration:
    \leanid{Matrix.HasSchmidtNumberLE.tensorMapId_tEta_posSemidef}.}
Composing that result with
(\ref{eq:wolf_reduction_criterion_schmidt_premise_formal_left}) and
(\ref{eq:wolf_reduction_criterion_schmidt_premise_formal_right}) yields
Equation~(3.18) of Ref.~\cite{Wolf2012Quantum} with its Schmidt-number
premise.\footnote{Formal declarations:
    \leanid{Matrix.reductionCriterion_left_of_hasSchmidtNumberLE} and
    \leanid{Matrix.reductionCriterion_right_of_hasSchmidtNumberLE}.}

The required positivity threshold is
\begin{align}
    1\leq n<D
        &\Longrightarrow T_n\ \text{is $n$-positive}.
    \label{eq:wolf_reduction_criterion_schmidt_premise_positivity_threshold}
\end{align}
This is formalized by
\leanid{Matrix.isNPositiveMap_tEta_iff} in
\doclink{QICLean/Channel/NPositivityChainStrict}; the reduction map is
$T_\eta$ at $\eta=n$ by \leanid{Matrix.reductionMap_eq_tEta}.

\section{The rectangular forward theorem}

The general forward step allows independent dimensions. Let
$T:M_d(\C)\to M_r(\C)$ be $n$-positive, and let
$\psi\in\C^d\otimes\C^k$ have Schmidt rank at most $n$. Write $\psi$ using a
normalized maximally entangled vector and a rectangular right-factor matrix
$X\in M_{d\times k}(\C)$ satisfying
\begin{align}
    \operatorname{rank}(X)
        &=\operatorname{SR}(\psi).
    \label{eq:wolf_reduction_criterion_schmidt_premise_rectangular_rank}
\end{align}
The ampliation of $\ketbra{\psi}{\psi}$ is a right-factor compression of the
rectangular Choi matrix of $T$. Its quadratic form pulls back to a Choi
quadratic form on a vector of Schmidt rank at most
$\operatorname{rank}(X)$. Hence
\begin{align}
    \operatorname{SR}(\psi)\leq n
        &\Longrightarrow
        (T\otimes\operatorname{id}_k)
        \!\left(\ketbra{\psi}{\psi}\right)\geq0.
    \label{eq:wolf_reduction_criterion_schmidt_premise_rectangular_pure_step}
\end{align}
Summing over a Schmidt-number decomposition gives
\begin{align}
    \operatorname{sn}(\rho)\leq n
        &\Longrightarrow
        (T\otimes\operatorname{id}_k)(\rho)\geq0.
    \label{eq:wolf_reduction_criterion_schmidt_premise_rectangular_mixed_step}
\end{align}
These results are formalized by
\leanid{Matrix.tensorMapId_posSemidef_of_hasSchmidtRankLE} and
\leanid{Matrix.HasSchmidtNumberLE.tensorMapId_posSemidef} in
\doclink{QICLean/Channel/SchmidtNumber}.

The older padding constructions
\leanid{Matrix.tensorMapId_posSemidef_of_hasSchmidtRankLE_general} and
\leanid{Matrix.HasSchmidtNumberLE.tensorMapId_posSemidef_general} remain as
compatibility declarations in
\doclink{QICLean/Channel/SchmidtNumberFactors}. The latter now delegates to the
direct rectangular theorem. Neither declaration is needed for the
source-facing positive-map characterization.

\section{Classification}

The bipartite-dimension scope restriction is resolved. The general implication
in
(\ref{eq:wolf_reduction_criterion_schmidt_premise_rectangular_mixed_step})
holds for every $n$-positive map with independent input, output, and bystander
dimensions through
\leanid{Matrix.HasSchmidtNumberLE.tensorMapId_posSemidef}.

The specialized reduction criterion is also dimension-general. The inequalities
\begin{align}
    n\Id\otimes\rho_2&\geq\rho,
    \label{eq:wolf_reduction_criterion_schmidt_premise_general_left_result}\\
    n\rho_1\otimes\Id&\geq\rho
    \label{eq:wolf_reduction_criterion_schmidt_premise_general_right_result}
\end{align}
are proved on a general space $\C^d\otimes\C^{d'}$ by
\leanid{Matrix.reductionCriterion_left_of_hasSchmidtNumberLE_general} and
\leanid{Matrix.reductionCriterion_right_of_hasSchmidtNumberLE_general}. These
theorems compose the rectangular forward result with
\leanid{Matrix.reductionCriterion_left} and
\leanid{Matrix.reductionCriterion_right}. The square declarations
\leanid{Matrix.reductionCriterion_left_of_hasSchmidtNumberLE} and
\leanid{Matrix.reductionCriterion_right_of_hasSchmidtNumberLE} remain the
specializations with $d=d'$.

The only residual scope concerns the positivity threshold of $T_n$: the left
inequality requires $1\leq n<d$, and the right inequality requires
$1\leq n<d'$, as in Equation~(3.11) of
Ref.~\cite{Wolf2012Quantum}. The complete-positivity endpoints $n=d$ and
$n=d'$, respectively, are documented in
\path{docs/paper-gaps/wolf_t_eta_top_index_scope.tex}. The dimension
restriction recorded in this note is fully eliminated.

\bibliographystyle{alpha}
\bibliography{references}

\end{document}
